\documentclass[11pt]{article}
\usepackage[margin=1in]{geometry}\usepackage{amsmath,amssymb,amsthm,microtype,hyperref}
\hypersetup{colorlinks=true,linkcolor=blue,urlcolor=blue}
\newtheorem{theorem}{Theorem}\newtheorem{lemma}[theorem]{Lemma}
\title{Exceptional cutoffs for the central Euler product\\at the boundary of short logarithmic intervals}
\author{Research note on a question in the ProveIt archive}\date{}
\begin{document}\maketitle
\begin{abstract}
For a fixed primitive nonprincipal Dirichlet character, the
deep Riemann hypothesis predicts an asymptotic for every sharp prime cutoff.
Using Dirichlet GRH and a frequency split in the explicit formula, we bound
the logarithmic measure of cutoffs with fixed canonical-log error greater
than $\epsilon$ in each interval $[e^U,e^{U+1}]$ by
$O_{\chi,\epsilon}(\exp[-\sqrt{2\pi\epsilon U}])$. The bound gives a
quantified response to the earlier ProveIt question about windows of
logarithmic length $\exp(-a\sqrt U)$. No independence or simplicity
assumption on the zeros is used.
\end{abstract}
\section{Statement}
Let $\chi$ be primitive and nonprincipal, $m=\operatorname{ord}_{1/2}L(s,\chi)$,
and $\delta=1$ if $\chi^2$ is principal, $0$ otherwise. The canonical
logarithm of the sharp product $P_x=\prod_{p\le x}(1-\chi(p)/\sqrt p)^{-1}$ is
\[
 \ell(x)=\sum_{p\le x}\sum_{k\ge1}\frac{\chi(p)^k}{k p^{k/2}} .
\]
There is a constant $C_\chi$ with
\[
 \exp C_\chi=2^{\delta/2}e^{-m\gamma_{\rm E}}L^{(m)}(1/2,\chi)/m!.
\]
Its imaginary part is fixed by the continued Euler logarithm from the right,
as in the proof below. Define $B(v)=\ell(e^v)+m\log v-C_\chi$.
\begin{theorem}\label{th:main}
For each fixed $\epsilon>0$ and every $U\ge2$,
\[
 \big|\{v\in[U,U+1]:|B(v)|>\epsilon\}\big|
 \ll_{\chi,\epsilon}\exp(-\sqrt{2\pi\epsilon U}).
\]
Thus the set of exceptional $x\ge e^2$ has finite measure for
$e^{b\sqrt{\log x}}\,dx/x$ whenever $b<\sqrt{2\pi\epsilon}$. In each
interval $[U,U+h]$ with $h=\exp(-a\sqrt U)$, $a>0$, the fraction with
$|B|>\epsilon$ is
$O_{\chi,\epsilon}(\exp[-(\sqrt{2\pi\epsilon}-a)\sqrt U])$.
\end{theorem}
At the boundary window scale the exceptional fraction therefore vanishes
for every fixed tolerance $\epsilon>a^2/(2\pi)$. This does not give
convergence for every tolerance in each such window at fixed $a$, or
a pointwise DRH theorem. For example, taking $a=1$ allows every tolerance
greater than $1/(2\pi)$ in the canonical logarithm. These estimates also
bound relative product discrepancies: $|B|\le\epsilon$ implies
$|(\log x)^m P_x e^{-C_\chi}-1|\le e^\epsilon-1$ at $x=e^v$.
\section{Reduction to an explicit formula}
Write $\Psi_\chi(x)=\sum_{n\le x}\chi(n)\Lambda(n)$ and
$E(v)=e^{-v/2}\Psi_\chi(e^v)$. We need the endpoint formula
\begin{equation}\label{eq:ep}
 B(v)=E(v)/v+O_\chi(1/v),\qquad v\ge2 .
\end{equation}
It is the fixed-center case of the archived endpoint calculation \cite{in};
here are the details. The von Mangoldt explicit formula on $v\ge1$, in
distributions, gives
\[
 E(v)=-\sum_{\rho=1/2+i\gamma}\frac{e^{i\gamma v}}{\rho}
           +O_\chi((1+v)e^{-v/2}).
\]
The last term is an ordinary function coming from the pole of the Perron
kernel and the trivial zeros. The $m$ central zeros give $-2m$. The
primitive of $E(v)+2m$ on $[1,\infty)$ is bounded, since integration
over a finite interval costs at most $2/|\rho\gamma|$ per nonzero
ordinate and the sum converges. Let
$H(v)=\sum_{2\le n\le e^v}\chi(n)\Lambda(n)/(n^{1/2}\log n)$.
Stieltjes summation gives an endpoint $E(v)/v$ and an integral with
integrand $E(w)(1/(2w)+1/w^2)$. The constant part yields
$-m\log v$ plus a constant and an $O(1/v)$ endpoint. Integration by
parts with the bounded primitive of the remaining part bounds its
integral from $v$ to infinity by $O_\chi(1/v)$. Consequently
$H(v)=c_\chi-m\log v+E(v)/v+O_\chi(1/v)$.
For $\sigma>0$, its Abel mean
$\sigma\int_0^\infty e^{-\sigma v}H(v)\,dv$ equals
$\log L(1/2+\sigma,\chi)$ by partial summation and continuation of the
Euler logarithm. In fact GRH gives $\Psi_\chi(x)\ll_\chi\sqrt{x}\log^2(2x)$,
ensuring convergence in this half-plane. The bounded primitive controls
the Abel mean of $(E(v)+2m)/v$ on $v\ge1$, and bounded lower intervals
contribute $o(1)$. Hence
$\log L(1/2+\sigma,\chi)=c_\chi+m\log\sigma+m\gamma_{\rm E}+o(1)$.

The ordinary prime cutoff adds
$\tfrac12\sum_{\sqrt x<p\le x}\chi^2(p)/p=\tfrac\delta2\log2+
O_\chi(1/\log x)$ by fixed-modulus prime number theory and partial
summation. Higher powers contribute $O(x^{-1/6})$ in absolute value:
for $p>x^{1/3}$ sum $O(p^{-3/2})$, and for smaller $p$ the first
omitted power contributes $O(x^{-1/2})$ per prime including its tail.
Thus $C_\chi=c_\chi+\tfrac\delta2\log2$ and \eqref{eq:ep} follows.
The canonical choice of $c_\chi$ is the constant in the displayed
expansion of $H$, before exponentiation.
\section{A tail bound uniform in the interval location}
\begin{lemma}\label{lem:tail}
For every real $\lambda\ge1$ and every $U\ge2$,
\[
 |\{v\in[U,U+1]: |E(v)|>\lambda\}|
       \ll_\chi e^{-\sqrt{2\pi\lambda}} .
\]
\end{lemma}
\begin{proof}
Count both signs with multiplicity. The Dirichlet zero count \cite{count}
gives $N(T)=(T/\pi)\log T+O_\chi(T)$ and
$N(t+1)-N(t)\ll_\chi\log(2+t)$. Partial summation gives, for $T\ge2$,
\[
 \sum_{|\gamma|\le T}|\rho|^{-1}
 \le \frac{(\log T)^2}{2\pi}+O_\chi(\log(2T)).
\]
Indeed one can replace $1/|\rho|$ by $1/|\gamma|$ above height one
at bounded total cost and integrate $N(t)/t^2$. The smooth remainder in
the explicit formula is uniformly bounded on $v\ge2$.

For the tail $S_T(v)=\sum_{|\gamma|>T} e^{i\gamma v}/\rho$,
\[
 \int_U^{U+1}|S_T(v)|^2\,dv\ll_\chi\frac{\log^2(2T)}{T}.
\]
For a finite subseries first majorize the unit interval by a constant
times a Gaussian centered at $U+1/2$. If $b_j$ is the sum of the absolute
coefficients in the bin $[j,j+1)$, Fourier transformation bounds the
integral by a constant times $\sum b_j^2$. Indeed the interaction of bins
$j,k$ has a bound $C e^{-c(\max\{|j-k|-1,0\})^2} b_jb_k$.
Now $b_j\ll_\chi\log(2+|j|)/(1+|j|)$ in the tail, which yields the
estimate by summing squares. The same bound makes the infinite series
converge in local $L^2$, and its distributional limit is the explicit
formula's tail.

For large $\lambda$ choose $\log T=\sqrt{2\pi\lambda}-D_\chi$, where the
fixed $D_\chi$ is large enough that the low-ordinate sum plus the smooth
remainder is at most $\lambda-\sqrt{\lambda}$. Expanding the square in
the low-ordinate bound proves this choice works. Chebyshev's inequality
then bounds the exceptional measure by
$O_\chi(\log^2(2T)/(\lambda T))=O_\chi(e^{-\sqrt{2\pi\lambda}})$.
Enlarge the constant for bounded $\lambda$.
\end{proof}
Apply the lemma at $\lambda=\epsilon U-O_\chi(1)$ using \eqref{eq:ep}.
This proves the first assertion of Theorem~\ref{th:main}. Sum over
integer $U$ with weight $e^{b\sqrt{U+1}}$ for the second, and divide
the unit-interval bound by $h$ for the third. No claim of optimality
for $2\pi$ is made; it comes from the two-sided zero count and the
cutoff used in this proof.

\section{Scope}
Sheth proves the central Euler-product prediction outside a set of finite
logarithmic measure for a broad automorphic class \cite{sheth}. The
archived Dirichlet analysis proves a local root-mean-square bound
$O_\chi(\log^2(2/h)/U)$ for $B$ and asks what happens at
$h=\exp(-a\sqrt U)$ \cite{in}. With Chebyshev's inequality that estimate
alone gives $O_\chi(a^4/\epsilon^2)$ there, without decay in $U$.
Theorem~\ref{th:main} gives decay for the tolerance range stated above.
The proof refines the familiar zero-truncation argument by matching its
cutoff to the error threshold and retaining the leading coefficient of the
zero count. The upstream Lean result supplies Dirichlet GRH; this tail
bound is an analytic deduction, with a classical proof. An independent
publication-priority claim would require a broader literature review.
\begin{thebibliography}{9}
\bibitem{in} V. Reshetnikov, ProveIt incoming collection,
\href{https://github.com/VladimirReshetnikov/ProveIt/blob/d31039a48d963692e5a5176f47b6603fc9e9d9cf/docs/incoming/dirichlet-grh-research.zip}{\texttt{dirichlet-grh-research.zip}},
\texttt{critical-line/critical-line.tex}, local mean square and research question 2.
\bibitem{count} M. A. Bennett, G. Martin, K. O'Bryant and A. Rechnitzer,
\href{https://doi.org/10.1090/mcom/3599}{Counting zeros of Dirichlet
$L$-functions}, Math. Comp. 90 (2021).
\bibitem{sheth} A. Sheth,
\href{https://doi.org/10.1017/S0305004125000234}{Euler products at the
centre and applications to Chebyshev's bias}, Math. Proc. Camb. Phil.
Soc. 179 (2025), 331--349, Theorems 3.2, 3.5 and 3.9.
\bibitem{conrad} K. Conrad,
\href{https://doi.org/10.4153/CJM-2005-012-6}{Partial Euler products
on the critical line}, Canad. J. Math. 57 (2005), 267--297.
\bibitem{lean} L. Sernau and collaborators,
\href{https://github.com/lsernau/grh-lean/tree/72fead102f663d5c7a407964a13dc0cf9d977207}{\emph{GRH for Dirichlet $L$-functions, in Lean 4}}.
\end{thebibliography}\end{document}
